% 2026 narrative draft; source mapping and open items are in planning/2026-draft-status.md.
% Shared presentation remains in shared/preamble.tex.
\firstname{David}
\familyname{Molik}
\title{Research, Scholarly and Creative Activity -- 2026}
\address{University Library}{Kansas State University}
\portfoliotitle

\section{A research agenda for reusable scientific information}
My research links computational genomics with the standards, software and publishing practices needed to make scientific information reusable. This is directly relevant to Libraries research support: questions about provenance, interoperability and sustainable access become both methodological problems and service-design requirements.

\subsection*{Publications and manuscripts}
I co-authored \textit{Toward standardization in arthropod and biodiversity genome projects} (\textit{Genetics}, 2026), extending my collaborative work on genomic reporting and interoperability. The article is a published output. In contrast, \textit{The Future is FAIR: A Community Framework for Enhanced Data Management and Data Sharing} is documented through a revised manuscript and reviewer response, and \textit{AgBioData Annual Meeting 2026: Toward Sustainable Agricultural Genomics Infrastructure} through a manuscript-preparation file. I report the latter two as work in progress, not accepted publications.\evidenceref{R1}

The cumulative CV records the earlier \textit{AgBioDatabase Finder}, BALROG-MON and genome-nomenclature articles. They provide context rather than new 2026 publication counts. In particular, BALROG-MON's February 19, 2025 publication precedes the February 25 K-State appointment recorded in the portfolio.

\subsection*{Awards, submissions and project development}
The following activity demonstrates development of a research program; requested budgets are not awarded funding.\evidenceref{G1}
\begin{itemize}
\item \textbf{ID3A seed funding, awarded:} a March 11 letter commits \$13,417 for ``Development of AI-ready digital ag research platform,'' starting March 12. Correspondence documents spending coordination; final expenditure reconciliation and reporting remain to be retained.
\item \textbf{Ag-DEx, submitted and not funded:} as PI/project director, I coordinated proposal components, collaboration materials, personnel and successive narrative drafts. The final package names Gwen Sibley as co-PD/co-PI. USDA received the application June 26 (2026-07893), requesting \$923,325 for January 2027--December 2029; the not-funded status was recorded August 27. Later response work addresses possible revisions rather than a completed resubmission.
\item \textbf{K-FORGE, submitted:} I served as PI with Nicholas Eggleston and Jeffrey Sheldon as co-PIs for ``Category II: K-FORGE: A Federated Open Research Gateway for Executable, Reproducible, and AI-Enabled Science.'' NSF's notification records July 28 submission, proposal 2638522. The request is \$9 million over 36 months, with a proposed February 2027 start. The proposal describes my intended leadership, coordination and governance responsibilities; no funding decision is established by the retained evidence.
\item \textbf{FFAR, submitted:} ``Unlocking Agricultural Research Data for Better Farming Decisions'' was submitted April 29, requesting \$449,912.20 for 2027--2029. This is a full application, not only institutional clearance; an award is not established.
\item \textbf{FAIROS, submitted then withdrawn:} NSF proposal 2624126, ``Disciplinary Improvements: AI-Ready Genome Notes: An Open Science Infrastructure for Rapid, Structured Genomic Reporting,'' was submitted April 8 and withdrawn June 8. Withdrawal is distinct from a sponsor decline.
\end{itemize}
SAS, Sunderland ATE/IUSE and OAK-AI development are described under Directed Service as partnership and project-scoping contributions. They are not added to the submitted-proposal or award totals.

\subsection*{Community research coordination and instruction}
As AgBioData co-PI and K-State subcontract PI, I contributed to the 2026 hybrid community meeting through budget, team and meeting coordination. The annual-report preview identifies these responsibilities but does not prove submission of that report. I also contributed to follow-up and meeting-report development. Survey calculations remain for a later review; I do not convert an AI-generated debrief into verified participation or learning metrics.\evidenceref{R2}

The May data-curation workshop records and follow-up correspondence document educational engagement and participant support. The schedule identifies May 1 and May 29 events; it does not establish my individual instructional hours or attendance totals. Earlier reported 2026 scholarly presentations remain in the evidence inventory for program/material reconciliation rather than being assigned uncertain conference editions here.\evidenceref{D1}

\subsection*{Research software and reproducibility}
I directed and reviewed software development extending my computational research. Barcode Simulator's August--September changes add distance-matrix and graph outputs, compressed/FASTQ output, reproducible run metadata and storage checks. These are software-maintenance and reproducibility contributions, not new biological findings; synthetic FASTQ qualities do not constitute an instrument-error model. PathwayByWire's August work includes pathway-completeness analysis, KEGG exports and an OpenAPI interface. Its conservation/evidence calculations should not be interpreted as measured expression or activity.\evidenceref{S1}

For HAIL, merged September changes implement three-valued bibliometric observations, quality signals, harvest coverage, Janeway enrichment and normalization replay. These are stronger evidence than architectural plans, but remain distinct from a validated production harvest. The institutional purpose and preliminary reporting outputs are described under Directed Service. Repository history includes collaborators and AI-assisted development; my contribution is research direction, design decisions and review, without claiming sole manual authorship of all code.\evidenceref{S1}\evidenceref{K1}
